\section*{Aktivitas Pembelajaran dan Praktikum}
\addcontentsline{toc}{section}{Aktivitas Pembelajaran dan Praktikum}

\begin{table}[H]
\centering
\caption{Rencana aktivitas Bab VI}
\label{tab:ch06-aktivitas}
\begin{tabularx}{\textwidth}{@{}llX@{}}
\toprule
\textbf{Sesi} & \textbf{Bentuk} & \textbf{Kegiatan} \\
\midrule
1 & Ceramah & DFA, NFA, dan epsilon closure \\
2 & Latihan & Menelusuri string pada DFA dan NFA \\
3 & Lokakarya & Subset construction, tabel konversi, dan minimalisasi \\
4 & Praktikum & P1-B: simulator NFA/DFA C dan uji ekivalensi \\
\bottomrule
\end{tabularx}
\end{table}

\subsection*{P1-B: Simulator NFA/DFA dan Tabel Konversi}
\textbf{Tujuan.} Membuat tabel konversi NFA ke DFA, tabel minimalisasi,
dan simulator NFA/DFA dalam C, lalu menguji bahwa keduanya menerima
bahasa token yang sama.

\textbf{Langkah kerja.}
\begin{enumerate}[label=\arabic*.,leftmargin=*]
  \item Untuk setiap kategori token pada \file{spec/tokens.md}, gambar
        NFA yang mengenali kategori tersebut.
  \item Lakukan subset construction dan isi tabel konversi: state DFA,
        himpunan state NFA, transisi setiap simbol, dan status
        penerima.
  \item Minimalisasi DFA dan isi tabel partisi: iterasi, partisi, serta
        alasan setiap pemisahan.
  \item Implementasikan simulator NFA dan DFA dalam C pada
        \file{nfa\_sim.c} dan \file{dfa\_sim.c}.
  \item Uji kedua simulator dengan 30 kasus uji dari P1-A.
  \item Pastikan NFA, DFA hasil konversi, dan DFA minimal menerima
        bahasa token yang sama.
  \item Bandingkan hasilnya dengan \file{fa\_tool.py}; catat perbedaan
        perilaku sebagai pembanding, bukan sebagai artefak utama.
\end{enumerate}

\textbf{Definisi selesai.}
\begin{checklist}
  \item Tabel konversi NFA ke DFA lengkap dan konsisten.
  \item Tabel minimalisasi memuat alasan setiap pemisahan.
  \item Simulator NFA dan DFA dalam C berjalan.
  \item 30 kasus uji menghasilkan keputusan yang sama pada ketiga
        automata.
  \item Perbedaan dengan \file{fa\_tool.py} terdokumentasi.
\end{checklist}

P1-B menjadi bagian Praktikum 1 (12\%) yang dilaporkan pada Bab VII.
